%&17pt
\pdfhorigin=5mm
\hoffset=12mm
\hsize=\pdfpagewidth \advance\hsize by-2\pdfhorigin \advance\hsize by-\hoffset
\pdfvorigin=12mm
\vsize=270mm
\output={\ifodd\pageno\else\hoffset=0pt\fi \plainoutput}

\input epsf

\nopagenumbers
\topglue 0pt
\vbox to0pt{\line{\vrule height\vsize \hfil\vrule}\vss}
\hrule
\kern2mm
\vbox to0pt{\centerline{\epsfxsize=10cm \epsfbox{upper.eps}}\vss}
\nointerlineskip
\line{\kern2mm \epsfxsize=30mm \epsfbox{upper-left.eps}\hfil
               \epsfxsize=30mm \epsfbox{upper-right.eps}\kern2mm}
\nointerlineskip
\vskip 0pt plus1filll
\vbox to0pt{\kern1cm
  \centerline{\font\F=omssbx17 scaled 7000 \F АСО}
  \kern2cm
  \centerline{\epsfxsize=10cm \epsfbox{center.eps}}
  \vss}
\nointerlineskip
\line{\kern2mm \epsfysize=15cm \epsfbox{left.eps}\hfil
               \epsfysize=15cm \epsfbox{right.eps}\kern2mm}
\nointerlineskip
\vskip 0pt plus1filll
\vbox to0pt{\vss\centerline{\font\F=omssbx14 \F Автоматическая
  Система Оповещения}}
\nointerlineskip
\line{\kern2mm \epsfxsize=30mm \epsfbox{lower-left.eps}\hfil
               \epsfxsize=30mm \epsfbox{lower-right.eps}\kern2mm}
\nointerlineskip
\vbox to0pt{\vss\centerline{\epsfxsize=10cm \epsfbox{lower.eps}}}
\kern2mm
\hrule
\vfill\eject
\line{}\vfill\eject

\topglue 2cm
\centerline{\font\bf=ombx17 at26pt \bf Содержание}
\kern2cm

\font\confont=omssdc10 at 22pt
\def\leaderfill{\leaders\hbox to 1em{\hss.\hss}\hfill}
\begingroup
\def\line#1{\hbox to\hsize{\hskip20mm #1\hskip20mm }}
\baselineskip=50pt

\line{\confont Запуск оповещения\leaderfill 1}
\line{\indent 1. Создание звукозаписи\leaderfill 1}
\line{\indent 2. Запуск оповещения\leaderfill 3}
\line{\indent 3. Печать отчёта\leaderfill 6}
\line{\indent 4. Очистка звукозаписи\leaderfill 7}

\line{\confont Редактирование списка оповещения\leaderfill 9}
\line{\indent Добавление абонента\leaderfill 9}
\line{\indent Удаление абонента\leaderfill 12}

\line{\confont Добавление абонента в адресную книгу\leaderfill 15}

\line{\confont Создание списка оповещения\leaderfill 19}

\endgroup
\vfill\eject
\line{}\vfill\eject

\font\secfont=omb10 at27pt
\font\subsecfont=ombx10 at20pt
\def\sec#1\par{\centerline{\secfont#1}\kern2cm}
\def\subsec#1\par{\leftline{\subsecfont#1}\medskip}

\def\rectangle#1#2#3#4{\vbox to0pt{\vss\hrule\kern-.3pt
  \hbox to#3{\vrule height#1 depth#2\hss#4\hss\vrule}%
  \kern-.3pt\hrule\kern-#2\kern-.1pt}}
\footline={\ifodd\pageno\hss\fi\vbox to0pt{\kern-5mm
  \rectangle{13pt}{2pt}{25pt}{\seventeenrm\folio}\vss}\ifodd\pageno\else\hss\fi}

\pageno=1

\sec Запуск оповещения

\subsec 1. Создание звукозаписи

\item{1)} Запустить программу записи.

$$\epsfbox{Запуск_оповещения/01-создать_звукозапись/01.eps}$$

\item{2)} Нажать микрофон, записать сообщение. В конце сообщения сказать
«ДЛЯ ПОДТВЕРЖДЕНИЯ НАЖМИТЕ ОДИН».

$$\epsfbox{Запуск_оповещения/01-создать_звукозапись/02.eps}$$

\item{3)} Остановить запись сообщения.

$$\epsfbox{Запуск_оповещения/01-создать_звукозапись/03.eps}$$

\item{4)} Прослушать последнюю запись (самая верхняя строка).
Если запись плохая, то записать новую.

$$\epsfbox{Запуск_оповещения/01-создать_звукозапись/04.eps}$$

\item{5)} Открыть папку «Звукозапись».

$$\epsfbox{Запуск_оповещения/01-создать_звукозапись/05.eps}$$

\item{6)} Щёлкнуть два раза на ogg с именем последней записи.

$$\epsfbox{Запуск_оповещения/01-создать_звукозапись/06.eps}$$

\item{7)} Вместо ogg должен появиться файл wav с таким же именем.

$$\epsfbox{Запуск_оповещения/01-создать_звукозапись/07.eps}$$

\subsec 2. Запуск оповещения

\item{1)} Запустить браузер либо выбрать открытый.

$$\epsfbox{Запуск_оповещения/02-запустить_оповещение/01.eps}$$

\item{2)} Запуск АСО (пропустить этот шаг если в первом пункте выбрали открытый браузер).

$$\epsfbox{Запуск_оповещения/02-запустить_оповещение/02.eps}$$

\item{3)} Вход в АСО (пропустить этот шаг если в первом пункте выбрали открытый браузер).

$$\epsfbox{Запуск_оповещения/02-запустить_оповещение/03.eps}$$

\item{4)} Выбрать в столбце слева «Группы оповещений».

$$\epsfbox{Запуск_оповещения/02-запустить_оповещение/04.eps}$$

\vfill\eject

\item{5)} Выбираем группу (щёлкаем по названию а не по квадратику слева от названия).

$$\epsfxsize=188mm \epsfbox{Запуск_оповещения/02-запустить_оповещение/05.eps}$$

\item{6)} Убедиться что была выбрана необходимая группа.

$$\epsfbox{Запуск_оповещения/02-запустить_оповещение/06.eps}$$

\item{7)} Загрузить файл звукозаписи, созданный как написано выше.

$$\epsfbox{Запуск_оповещения/02-запустить_оповещение/07.eps}$$

\item{8)} Выбрать файл.

$$\epsfbox{Запуск_оповещения/02-запустить_оповещение/08.eps}$$

\vfill\eject

\item{9)} Подтвердить выбор.

$$\epsfbox{Запуск_оповещения/02-запустить_оповещение/09.eps}$$

\item{10)} Сохранить.

$$\epsfbox{Запуск_оповещения/02-запустить_оповещение/10.eps}$$

\item{11)} Запустить оповещение.

$$\epsfbox{Запуск_оповещения/02-запустить_оповещение/11.eps}$$

\item{12)} Подтвердить запуск.

$$\epsfbox{Запуск_оповещения/02-запустить_оповещение/12.eps}$$

\vfill\eject

\subsec 3. Печать отчёта

\item{1)} Выбрать в столбце слева «Сеансы оповещений».

$$\epsfbox{Запуск_оповещения/03-печать_отчета/01.eps}$$

\item{2)} Сформировать отчёт.

$$\epsfbox{Запуск_оповещения/03-печать_отчета/02.eps}$$

\item{3)} Скачать PDF.

$$\epsfbox{Запуск_оповещения/03-печать_отчета/03.eps}$$

\item{4)} Открыть PDF.

$$\epsfbox{Запуск_оповещения/03-печать_отчета/04.eps}$$

\vfill\eject

\item{5)} Щёлкнуть «Файл» и в появившемся меню выбрать «Печать».

$$\epsfbox{Запуск_оповещения/03-печать_отчета/05.eps}$$

\item{6)} Распечатать отчёт.

$$\epsfbox{Запуск_оповещения/03-печать_отчета/06.eps}$$

\subsec 4. Очистка звукозаписи

\item{1)} Закрыть программу звукозаписи.

$$\epsfbox{Запуск_оповещения/04-очистить_звукозапись/01.eps}$$

\vfill\eject

\item{2)} Открыть папку «Звукозапись».

$$\epsfbox{Запуск_оповещения/04-очистить_звукозапись/02.eps}$$

\item{3)} Щёлкнуть правой кнопкой на каждом файле и в появившемся меню выбрать «Удалить».

$$\epsfbox{Запуск_оповещения/04-очистить_звукозапись/03.eps}$$

\vfill\eject

\sec Редактирование списка оповещения

\subsec Добавление абонента

\item{1)} Запустить браузер либо выбрать открытый.

$$\epsfbox{Редактирование_списка_оповещения/Добавление_абонента/01.eps}$$

\item{2)} Запуск АСО (пропустить этот шаг если в первом пункте выбрали открытый браузер).

$$\epsfbox{Редактирование_списка_оповещения/Добавление_абонента/02.eps}$$

\item{3)} Вход в АСО (пропустить этот шаг если в первом пункте выбрали открытый браузер).

$$\epsfbox{Редактирование_списка_оповещения/Добавление_абонента/03.eps}$$

\vfill\eject

\item{4)} Выбрать в столбце слева «Группы оповещений».

$$\epsfbox{Редактирование_списка_оповещения/Добавление_абонента/04.eps}$$

\item{5)} Выбрать список (нажать на само слово а не на строку).

$$\epsfxsize=188mm \epsfbox{Редактирование_списка_оповещения/Добавление_абонента/05.eps}$$

\item{6)} Добавить абонента.

$$\epsfbox{Редактирование_списка_оповещения/Добавление_абонента/06.eps}$$

\vfill\eject

\item{7)} Щёлкнуть в поле ввода.

$$\epsfbox{Редактирование_списка_оповещения/Добавление_абонента/07.eps}$$

\item{8)} Выбрать абонента (если абонента нет в выпадающем списке, то добавить
по разделу «Добавление абонента в адресную книгу»).

$$\epsfbox{Редактирование_списка_оповещения/Добавление_абонента/08.eps}$$

\item{9)} Выбрать номер.

$$\epsfbox{Редактирование_списка_оповещения/Добавление_абонента/09.eps}$$

\vfill\eject

\item{10)} Выбрать дополнительный номер.

$$\epsfbox{Редактирование_списка_оповещения/Добавление_абонента/10.eps}$$

\item{11)} Сохранить.

$$\epsfbox{Редактирование_списка_оповещения/Добавление_абонента/11.eps}$$

\subsec Удаление абонента

\item{1)} Запустить браузер либо выбрать открытый.

$$\epsfbox{Редактирование_списка_оповещения/Удаление_абонента/01.eps}$$

\item{2)} Запуск АСО (пропустить этот шаг если в первом пункте выбрали открытый браузер).

$$\epsfbox{Редактирование_списка_оповещения/Удаление_абонента/02.eps}$$

\item{3)} Вход в АСО (пропустить этот шаг если в первом пункте выбрали открытый браузер).

$$\epsfbox{Редактирование_списка_оповещения/Удаление_абонента/03.eps}$$

\item{4)} Выбрать в столбце слева «Группы оповещений».

$$\epsfbox{Редактирование_списка_оповещения/Удаление_абонента/04.eps}$$

\item{5)} Выбрать список (нажать на само слово а не на строку).

$$\epsfxsize=188mm \epsfbox{Редактирование_списка_оповещения/Удаление_абонента/05.eps}$$

\vfill\eject

\item{6)} Выбрать одного или нескольких абонентов.

$$\epsfbox{Редактирование_списка_оповещения/Удаление_абонента/06.eps}$$

\item{7)} Удалить выбранное.

$$\epsfbox{Редактирование_списка_оповещения/Удаление_абонента/07.eps}$$

\vfill\eject

\sec Добавление абонента в адресную книгу

\item{1)} Запустить браузер либо выбрать открытый.

$$\epsfbox{Добавление_абонента_в_адресную_книгу/01.eps}$$

\item{2)} Запуск АСО (пропустить этот шаг если в первом пункте выбрали открытый браузер).

$$\epsfbox{Добавление_абонента_в_адресную_книгу/02.eps}$$

\item{3)} Вход в АСО (пропустить этот шаг если в первом пункте выбрали открытый браузер).

$$\epsfbox{Добавление_абонента_в_адресную_книгу/03.eps}$$

\vfill\eject

\item{4)} Выбрать в столбце слева «Контакты».

$$\epsfbox{Добавление_абонента_в_адресную_книгу/04.eps}$$

\item{5)} Добавить контакт.

$$\epsfbox{Добавление_абонента_в_адресную_книгу/05.eps}$$

\vfill\eject

\item{6)} Заполнить ФИО и один или более номеров телефонов.

$$\epsfxsize=188mm \epsfbox{Добавление_абонента_в_адресную_книгу/06.eps}$$

\vfill\eject

\item{7)} Сохранить.

$$\epsfbox{Добавление_абонента_в_адресную_книгу/07.eps}$$

\vfill\eject

\sec Создание списка оповещения

\item{0)} Создать звукозапись как написано в подразделе \hfil\break
\hbox{«1. Создание звукозаписи»}
раздела \hbox{«Запуск оповещения»} \hfil\break (стр.~1).

\vskip1.5cm

\item{1)} Запустить браузер либо выбрать открытый.

$$\epsfbox{Создание_списка_оповещения/01.eps}$$

\item{2)} Запуск АСО (пропустить этот шаг если в первом пункте выбрали открытый браузер).

$$\epsfbox{Создание_списка_оповещения/02.eps}$$

\item{3)} Вход в АСО (пропустить этот шаг если в первом пункте выбрали открытый браузер).

$$\epsfbox{Создание_списка_оповещения/03.eps}$$

\vfill\eject

\item{4)} Выбрать в столбце слева «Списки контактов».

$$\epsfbox{Создание_списка_оповещения/04.eps}$$

\item{5)} Добавить список.

$$\epsfbox{Создание_списка_оповещения/05.eps}$$

\item{6)} Задать название списка. Начать со слова «Список».

$$\epsfbox{Создание_списка_оповещения/06.eps}$$

\item{7)} Сохранить изменения.

$$\epsfbox{Создание_списка_оповещения/07.eps}$$

\item{8)} Выбрать в столбце слева «Группы оповещений».

$$\epsfbox{Создание_списка_оповещения/08.eps}$$

\item{9)} Создать новую группу.

$$\epsfbox{Создание_списка_оповещения/09.eps}$$

\item{10)} Задать название группы такое же как название списка,
но заменить «Список» на «Группа» в начале названия.

$$\epsfbox{Создание_списка_оповещения/10.eps}$$

\item{11)} Загрузить файл звукозаписи, созданный на шаге ``0).''

$$\epsfbox{Создание_списка_оповещения/11.eps}$$

\vfill\eject

\item{12)} Выбрать файл.

$$\epsfbox{Создание_списка_оповещения/12.eps}$$

\item{13)} Подтвердить выбор файла.

$$\epsfbox{Создание_списка_оповещения/13.eps}$$

\item{14)} Задать номер 3055 (именно такой, всегда).

$$\epsfbox{Создание_списка_оповещения/14.eps}$$

\item{15)} Задать количество одновременных звонков: 10.

$$\epsfbox{Создание_списка_оповещения/15.eps}$$

\item{16)} Задать количество попыток: 2.

$$\epsfbox{Создание_списка_оповещения/16.eps}$$

\item{17)} Задать время ожидания ответа: 30 сек.

$$\epsfbox{Создание_списка_оповещения/17.eps}$$

\item{18)} Задать интервал между попытками: 30 сек.

$$\epsfbox{Создание_списка_оповещения/18.eps}$$

\item{19)} Задать код подтверждения: 1.

$$\epsfbox{Создание_списка_оповещения/19.eps}$$

\item{20)} Добавить одноимённый список в группу.

$$\epsfbox{Создание_списка_оповещения/20.eps}$$

\item{21)} Выбрать список.

$$\epsfbox{Создание_списка_оповещения/21.eps}$$

\item{22)} Добавить канал телефонной станции.

$$\epsfbox{Создание_списка_оповещения/22.eps}$$

\item{23)} Выбрать канал.

$$\epsfbox{Создание_списка_оповещения/23.eps}$$

\vfill\eject

\item{24)} Сохранить изменения.

$$\epsfbox{Создание_списка_оповещения/24.eps}$$

\item{25)} Завершить работу со звукозаписью как написано в подразделе \hfil\break
\hbox{«4. Очистка звукозаписи»} раздела \hbox{«Запуск оповещения»} \hfil\break (стр.~7).

\bye
